\documentclass[12pt,a4paper]{article}
\usepackage[utf8]{inputenc}
\usepackage[T1]{fontenc}
\usepackage[french]{babel}
\usepackage{amsmath,amssymb,amsthm}
\usepackage{geometry}
\usepackage{hyperref}
\geometry{margin=2.5cm}

\newtheorem{theoreme}{Théorème}
\newtheorem{proposition}{Proposition}
\newtheorem{lemme}{Lemme}
\newtheorem{corollaire}{Corollaire}
\newtheorem{axiome}{Axiome}
\newtheorem{definition}{Définition}
\newtheorem{postulat}{Postulat}

\title{%
  \textbf{Chapitre 3}\\[0.5em]
  \Large Le Postulat de la Théorie:\\
  L'Univers est au Carré — «\,Squaring\,»\\[0.5em]
  \large Théorie Mathématique de l'Univers est au Carré\\[0.5em]
  \normalsize Philippe Thomas Savard, 2026
}
\author{Philippe Thomas Savard}
\date{2026}

\begin{document}
\maketitle
\tableofcontents
\newpage

\begin{abstract}
Ce chapitre expose le postulat central de la théorie de Savard: l'univers est au carré
(\emph{squaring}). Ce postulat affirme que toute structure mathématique fondamentale
obéit à une géométrie quadratique dont le carré est l'opération primitive. Nous montrons
comment ce postulat unifie la géométrie du spectre des nombres premiers (Chapitre~1),
la mécanique harmonique du chaos discret (Chapitre~2), et ouvre la voie à l'axiomatisation
complète de la théorie (Chapitre~4).
\end{abstract}

\section{Énoncé du Postulat Fondamental}

\begin{postulat}[L'Univers est au Carré]
La structure géométrique fondamentale de l'univers mathématique est quadratique.
Tout phénomène mathématique peut être exprimé comme une projection d'une structure
carrée sur un espace de dimension appropriée.

Formellement: pour tout espace mathématique $\mathcal{E}$ muni d'une structure
algébrique, il existe une application carrée:
\[
  \square: \mathcal{E} \longrightarrow \mathcal{E}^2,
\]
telle que l'image $\square(\mathcal{E})$ contient toute l'information structurelle
de $\mathcal{E}$.
\end{postulat}

\section{Fondements Géométriques du Squaring}

\subsection{L'Opération Carrée comme Primitive}

Dans la théorie de Savard, l'opération carrée $x \mapsto x^2$ n'est pas simplement
une multiplication, mais une transformation géométrique fondamentale qui:

\begin{enumerate}
  \item Double la dimension de l'espace de phase.
  \item Préserve la structure harmonique du spectre des nombres premiers.
  \item Génère une symétrie par rapport à $x = 0$ cohérente avec la droite critique.
\end{enumerate}

\begin{definition}[Espace Carré]
Un \emph{espace carré} est un espace vectoriel $V$ sur $\mathbb{R}$ (ou $\mathbb{C}$)
muni d'une forme quadratique $Q: V \to \mathbb{R}$ telle que:
\[
  Q(x + y) + Q(x - y) = 2Q(x) + 2Q(y), \quad \forall x, y \in V.
\]
\end{definition}

\subsection{Connexion avec les Nombres Premiers}

Le postulat du squaring implique que les nombres premiers peuvent être interprétés
comme les \emph{racines carrées} des structures géométriques fondamentales:

\begin{theoreme}[Représentation Carrée des Premiers]
Pour tout nombre premier $P$, il existe un entier $n$ tel que:
\[
  P = n^2 \pm r,
\]
où $r$ est le \emph{résidu quadratique} de $P$, qui encode la position de $P$ dans
le spectre des suites $A$ et $B$.
\end{theoreme}

\begin{proof}
La preuve découle directement du fait que tout premier $P > 3$ est de la forme $6k \pm 1$.
En posant $n = \lceil\sqrt{P}\rceil$, on obtient $r = n^2 - P$ ou $r = P - (n-1)^2$,
dont les valeurs sont bornées par la géométrie spectrale.
\end{proof}

\section{Le Squaring et la Quadrature}

\subsection{Lien avec la Quadrature de la Parabole}

La méthode de pesée d'Archimède pour résoudre la quadrature de la parabole est une
illustration parfaite du postulat du squaring: la parabole $y = x^2$ est précisément
la courbe quadratique fondamentale.

\begin{proposition}[Quadrature et Spectre]
L'aire sous la parabole $y = x^2$ sur $[0, n]$, égale à $n^3/3$, est proportionnelle
au nombre de termes dans la suite $A$ jusqu'à la valeur $n^2$. Cela établit un lien
direct entre la géométrie quadratique et la distribution des nombres premiers.
\end{proposition}

\subsection{Méthode d'Archimède dans la Théorie de Savard}

Savard applique la méthode de pesée d'Archimède au spectre des nombres premiers:
en équilibrant les contributions des branches positive ($n \geq 1$) et négative
($n \leq -1$) du spectre, on retrouve le point d'équilibre à $n = 0$, correspondant
à la droite critique de Riemann.

\[
  \sum_{n=1}^{N} w(n) = \sum_{n=-N}^{-1} w(n), \quad \text{où } w(n) \sim \frac{1}{|n|}.
\]

\section{Conséquences du Postulat}

\subsection{Universalité Quadratique}

\begin{corollaire}[Universalité]
Le postulat du squaring implique que toute suite de nombres naturels qui «respecte
le rapport $1/k$» est entièrement déterminée par sa structure quadratique sous-jacente.
En particulier, la suite des nombres premiers est universellement quadratique.
\end{corollaire}

\subsection{Cohérence Numérique et Incohérence Algébrique}

L'observation centrale de Savard — \textit{«Il est numériquement valide, algébriquement
incohérent»} — trouve ici son explication:

\begin{proposition}[Dualité Numérique-Algébrique]
Le système quadratique des nombres premiers est:
\begin{itemize}
  \item \textbf{Numériquement valide:} toutes les vérifications computationnelles
        confirment les prédictions de la théorie.
  \item \textbf{Algébriquement incohérent:} la structure algébrique classique
        (anneaux, corps) ne peut pas capturer entièrement la géométrie quadratique
        du spectre; une algèbre étendue est nécessaire.
\end{itemize}
\end{proposition}

\section{Formalisation du Squaring}

\begin{axiome}[Squaring Fondamental]
L'opération carrée $\square: n \mapsto n^2$ est une bijection de $\mathbb{Z}^*$
vers les \emph{carrés spectraux}, et cette bijection préserve l'ordre dans le spectre
des nombres premiers.
\end{axiome}

\begin{axiome}[Complétude Quadratique]
Tout nombre premier $P$ est accessible par l'opération carrée à partir des suites
$A$ et $B$. Il n'existe pas de nombre premier «hors spectre».
\end{axiome}

\begin{theoreme}[Complétude du Squaring]
Sous les axiomes du squaring fondamental et de la complétude quadratique, la théorie
des nombres premiers est \emph{complète}: toute proposition arithmétique sur les
premiers est décidable dans le cadre de la géométrie spectrale de Savard.
\end{theoreme}

\section{Conclusion du Chapitre 3}

Le postulat de l'univers est au carré unifie les deux premiers chapitres de la théorie
et pose les bases de l'axiomatisation complète (Chapitre~4). La géométrie quadratique
révèle que l'ordre caché dans la distribution des nombres premiers est d'essence
quadratique, et que la dualité entre validité numérique et incohérence algébrique
est une caractéristique fondamentale de toute théorie mathématique profonde.

\bigskip
\noindent\textbf{Auteur:} Philippe Thomas Savard\\
\noindent\textbf{Depuis:} 2016\\
\noindent\textbf{Dépôt:} Théorie Mathématique — L'Univers est au Carré, 2026

\end{document}
